% Estilo común de diapositivas de sesión (PDF / Beamer 16:9). No usar PPTX.
\usepackage[T1]{fontenc}
\usepackage[spanish,es-tabla]{babel}
\usepackage{booktabs}
\usepackage{array}
\usepackage{colortbl}
\usepackage{xcolor}
\usepackage{hyperref}

\usetheme{Madrid}
\usecolortheme{whale}
\usefonttheme{professionalfonts}
\usepackage[scaled=0.92]{helvet}
\renewcommand{\familydefault}{\sfdefault}

\definecolor{uninavy}{RGB}{27,54,93}
\definecolor{uniaccent}{RGB}{232,163,23}
\definecolor{unimuted}{RGB}{99,110,114}
\definecolor{tablehead}{RGB}{27,54,93}
\definecolor{tablealt}{RGB}{238,242,247}

\setbeamercolor{structure}{fg=uninavy}
\setbeamercolor{frametitle}{bg=uninavy,fg=white}
\setbeamercolor{title}{fg=white,bg=uninavy}
\setbeamercolor{footline}{fg=unimuted,bg=white}
\setbeamertemplate{navigation symbols}{}
\setbeamersize{text margin left=10mm, text margin right=10mm}
\setbeamerfont{frametitle}{size=\large,series=\bfseries}

\setbeamertemplate{footline}{%
  \leavevmode%
  \hbox{%
    \begin{beamercolorbox}[wd=.78\paperwidth,ht=2.6ex,dp=1.2ex,leftskip=1em]{footline}%
      \insertshortauthor{}%
    \end{beamercolorbox}%
    \begin{beamercolorbox}[wd=.22\paperwidth,ht=2.6ex,dp=1.2ex,rightskip=1em]{footline}%
      \hfill\insertframenumber{}/\inserttotalframenumber\hspace{1em}%
    \end{beamercolorbox}%
  }%
}

\newcommand{\tabhead}[1]{\textcolor{white}{\textbf{#1}}}
\newcolumntype{L}[1]{>{\raggedright\arraybackslash}p{#1}}
